\documentclass[11pt,a4paper]{article}
% Fix listings UTF-8 compatibility
\usepackage{textcomp}

% ── Packages ──────────────────────────────────────────────────────────────
\usepackage[margin=1in]{geometry}
\usepackage[T1]{fontenc}
\usepackage[utf8]{inputenc}
\usepackage{microtype}
\usepackage{amsmath}
\usepackage{amssymb}
\usepackage{graphicx}
\usepackage{hyperref}
\usepackage{xcolor}
\usepackage{listings}
\usepackage{booktabs}
\usepackage{enumitem}
\usepackage{titlesec}
\usepackage{parskip}
\usepackage{fancyhdr}
\usepackage{setspace}

% ── Colours ───────────────────────────────────────────────────────────────
\definecolor{codebg}{RGB}{245,247,250}
\definecolor{keyblue}{RGB}{0,90,180}
\definecolor{commentgray}{RGB}{100,110,120}
\definecolor{stringgreen}{RGB}{0,130,60}
\definecolor{accent}{RGB}{30,80,160}

% ── Hyperlinks ────────────────────────────────────────────────────────────
\hypersetup{
  colorlinks=true,
  linkcolor=accent,
  urlcolor=accent,
  citecolor=accent
}

% ── Code listings style ───────────────────────────────────────────────────
\lstdefinestyle{pylucene}{
  language=Python,
  backgroundcolor=\color{codebg},
  basicstyle=\ttfamily\footnotesize,
  keywordstyle=\color{keyblue}\bfseries,
  commentstyle=\color{commentgray}\itshape,
  stringstyle=\color{stringgreen},
  numberstyle=\tiny\color{commentgray},
  numbers=left,
  stepnumber=1,
  numbersep=6pt,
  frame=single,
  framerule=0.4pt,
  rulecolor=\color{gray!40},
  breaklines=true,
  breakatwhitespace=true,
  showstringspaces=false,
  tabsize=4,
  captionpos=b,
  xleftmargin=12pt,
  xrightmargin=4pt,
  aboveskip=8pt,
  belowskip=8pt,
}
\lstset{style=pylucene}

% ── Section formatting ────────────────────────────────────────────────────
\titleformat{\section}{\large\bfseries\color{accent}}{}{0em}{}[\titlerule]
\titleformat{\subsection}{\normalsize\bfseries\color{accent!80}}{}{0em}{}
\titlespacing{\section}{0pt}{14pt}{6pt}
\titlespacing{\subsection}{0pt}{10pt}{4pt}

% ── Header / Footer ───────────────────────────────────────────────────────
\pagestyle{fancy}
\fancyhf{}
\rhead{\small Information Retrieval with PyLucene}
\lhead{\small \today}
\cfoot{\thepage}
\renewcommand{\headrulewidth}{0.4pt}

% ─────────────────────────────────────────────────────────────────────────
\begin{document}

% ── Title Page ────────────────────────────────────────────────────────────
\begin{titlepage}
  \centering
  \vspace*{1.5cm}
  {\Huge\bfseries\color{accent} Information Retrieval \\ with PyLucene \par}
  \vspace{0.6cm}
  \rule{0.6\linewidth}{1.2pt}
  \vspace{0.5cm}

  {\large Study Oriented Project (SOP) Final Report \par}
  \vspace{0.4cm}
  {\normalsize Covering IR Foundations, Ranking Models, and\\
               PyLucene Indexing \& Retrieval in Depth\par}

  \vspace{1.5cm}
  {\large \textbf{Submitted by:} \par}
  \vspace{0.2cm}
  {\large Aayush Vij \& Vandan \par}
  
  \vspace{1.5cm}
  {\large \textbf{Submitted to:} \par}
  \vspace{0.2cm}
  {\large Prof. Abhisek Chakrabarty \par}

  \vspace{1.5cm}
  \begin{tabular}{ll}
    \textbf{Repository} &
    \href{https://github.com/aayushvij04/IR}{\texttt{github.com/aayushvij04/IR}} \\
    \textbf{Date}       & \today \\
  \end{tabular}

  \vfill
  \small\textit{Source code and this report are available in the linked repository.}
\end{titlepage}

% ─────────────────────────────────────────────────────────────────────────
\section{Information Retrieval — Foundations}
% ─────────────────────────────────────────────────────────────────────────

\textbf{Information Retrieval (IR)} is the science of finding documents relevant to an information need from large collections. Unlike database retrieval, IR operates on unstructured text and copes with ambiguity and redundancy. A classical IR system has three pillars:
\begin{enumerate}[leftmargin=*, label=\arabic*., nosep]
  \item \textbf{Corpus and Documents.} A collection of textual units over which retrieval is performed.
  \item \textbf{Query Processing.} Translating an information need into Boolean expressions, vectors, or probabilities.
  \item \textbf{Ranking Model.} A mathematical model that assigns a relevance score to each document.
\end{enumerate}

\subsection{The Inverted Index}
The \emph{inverted index} is the foundational data structure. For every vocabulary term $t$, it stores a \emph{postings list}—a sorted list of document IDs:
\vspace{-0.5em}
\begin{equation}
  \text{Index}: t \longrightarrow \langle df_t,\;[d_1, d_2, \ldots, d_{df_t}] \rangle
\end{equation}
\vspace{-0.5em}

Boolean AND queries are answered in $O(|p_1| + |p_2|)$ via a linear merge of sorted lists, accelerated with \emph{skip pointers} every $\lfloor\sqrt{n}\rfloor$ positions. Construction at scale uses \textbf{BSBI} (Block Sort-Based Indexing) or \textbf{SPIMI} (Single-Pass In-Memory Indexing). Postings are compressed with Variable-Byte or Gamma encoding to reduce footprint, motivated by Zipf's Law and Heaps' Law ($M \approx k\,n^b$).

\subsection{Ranking Models}

\textbf{Vector Space Model (VSM).} Documents and queries are TF-IDF weighted vectors in $\mathbb{R}^{|V|}$. Relevance is the \emph{cosine similarity}:
\vspace{-0.5em}
\begin{equation}
  \text{sim}(q, d) = \frac{\vec{q} \cdot \vec{d}}{|\vec{q}|\;|\vec{d}|} \qquad \text{where}\quad w_{t,d} = \bigl(1 + \log \text{tf}_{t,d}\bigr)\cdot\log\frac{N}{df_t}
\end{equation}
\vspace{-0.5em}

\textbf{BM25 (Okapi).} The probabilistic ranking function dominating modern search engines:
\vspace{-0.5em}
\begin{equation}
  \text{BM25}(q,d) = \sum_{t\in q} \text{IDF}(t) \cdot \frac{f_{t,d}(k_1+1)}{f_{t,d}+k_1(1-b+b\frac{|d|}{\text{avgdl}})}
\end{equation}
\vspace{-0.5em}

Typical parameters are $k_1 \in [1.2, 2.0]$ and $b = 0.75$. Extensions include \textbf{BM25F} (field-weighted) and \textbf{BM25+} (lower-bound on TF saturation).

\textbf{Language Models (LM).} Query likelihood is $P(q\mid d) = \prod_{t\in q} P(t\mid d)$. The zero-frequency problem is solved with smoothing methods such as \emph{Jelinek-Mercer}: $P_\lambda(t\mid d) = \lambda P(t\mid d)_\text{MLE} + (1-\lambda)P(t\mid C)$, or \emph{Dirichlet}: $P_\mu(t\mid d) = \frac{f_{t,d}+\mu P(t\mid C)}{|d|+\mu}$.

\subsection{Evaluation}
Retrieval quality is measured by \textbf{MAP} (Mean Average Precision), \textbf{nDCG@K} (Normalized Discounted Cumulative Gain), and \textbf{MRR} (Mean Reciprocal Rank). Evaluation often follows the \textbf{TREC} paradigm using pooled judgements and \texttt{trec\_eval}. Inter-annotator agreement is quantified by Cohen's $\kappa$.

% ─────────────────────────────────────────────────────────────────────────
\newpage
\section{PyLucene — Architecture and Setup}
% ─────────────────────────────────────────────────────────────────────────

\textbf{Apache Lucene} is a high-performance, open-source full-text search
library written in Java. \textbf{PyLucene} exposes the full Lucene API to Python
via JCC, a C++ code generator that bridges the JVM and CPython runtimes.
This means every Java class in Lucene 9.x is directly accessible as a Python
class—with identical semantics and performance.

\subsection{Initializing the JVM}

Every PyLucene session must start the Java Virtual Machine exactly once:

\begin{lstlisting}[caption={JVM initialization}]
import lucene
lucene.initVM(vmargs=['-Djava.awt.headless=true'])
\end{lstlisting}

\subsection{Storage Backends}

\begin{lstlisting}[caption={On-disk vs in-memory index}]
from org.apache.lucene.store import FSDirectory, ByteBuffersDirectory
from java.nio.file import Paths

# Persistent on-disk index
directory = FSDirectory.open(Paths.get("/path/to/index"))

# Ephemeral in-memory index (testing / unit tests)
directory = ByteBuffersDirectory()
\end{lstlisting}

% ─────────────────────────────────────────────────────────────────────────
\section{Indexing with PyLucene}
% ─────────────────────────────────────────────────────────────────────────

\subsection{IndexWriter and Analyzers}

\begin{lstlisting}[caption={Creating an IndexWriter with StandardAnalyzer}]
from org.apache.lucene.analysis.standard import StandardAnalyzer
from org.apache.lucene.index import IndexWriter, IndexWriterConfig

analyzer = StandardAnalyzer()          # lowercase + stop-word removal
config   = IndexWriterConfig(analyzer)
config.setOpenMode(IndexWriterConfig.OpenMode.CREATE)
writer   = IndexWriter(directory, config)
\end{lstlisting}

Available analyzers and their behaviours are summarized in Table~\ref{tab:analyzers}.

\begin{table}[h]
\centering
\small
\begin{tabular}{@{}lll@{}}
\toprule
\textbf{Analyzer} & \textbf{Tokenization} & \textbf{Typical Use} \\
\midrule
\texttt{StandardAnalyzer}   & Unicode tokenizer + lowercase + stop words & General English text \\
\texttt{WhitespaceAnalyzer} & Split on whitespace, preserves case         & Code, identifiers    \\
\texttt{KeywordAnalyzer}    & Entire field = one token                    & IDs, enums, genres   \\
\texttt{EnglishAnalyzer}    & StandardAnalyzer + Porter stemming          & IR experiments       \\
\texttt{CJKAnalyzer}        & Bi-gram tokenization                        & Chinese/Japanese/Korean \\
\bottomrule
\end{tabular}
\caption{PyLucene Analyzer options}
\label{tab:analyzers}
\end{table}

\subsection{Document and Field Types}

A Lucene \texttt{Document} is a flat set of \texttt{Field} objects.
Field type controls \emph{whether} and \emph{how} a value is indexed:

\begin{lstlisting}[caption={Field types in a movie document (from \texttt{src/pylucene/indexing.py})}]
from org.apache.lucene.document import (
    Document, TextField, StringField,
    StoredField, IntPoint, FloatPoint, Field)

def movie_to_document(movie: dict) -> Document:
    doc = Document()
    # Analyzed full-text - tokenized, lowercased, stop-word removed
    doc.add(TextField("title", movie["title"], Field.Store.YES))
    doc.add(TextField("plot",  movie["plot"],  Field.Store.YES))
    # Exact keyword - single token, not analyzed
    doc.add(StringField("genre", movie["genre"], Field.Store.YES))
    # Stored only - retrieve at search time, not searchable
    doc.add(StoredField("year_display", str(movie["year"])))
    # Numeric range-queryable fields
    doc.add(IntPoint  ("year",   movie["year"]))
    doc.add(FloatPoint("rating", movie["rating"]))
    doc.add(StoredField("year",   movie["year"]))
    doc.add(StoredField("rating", movie["rating"]))
    return doc
\end{lstlisting}

\begin{table}[h]
\centering
\small
\begin{tabular}{@{}llll@{}}
\toprule
\textbf{Field Type} & \textbf{Indexed} & \textbf{Analyzed} & \textbf{Stored} \\
\midrule
\texttt{TextField}   & Yes & Yes & Optional \\
\texttt{StringField} & Yes & No  & Optional \\
\texttt{StoredField} & No  & No  & Yes      \\
\texttt{IntPoint}    & Yes (BKD tree) & No & No \\
\texttt{FloatPoint}  & Yes (BKD tree) & No & No \\
\bottomrule
\end{tabular}
\caption{Lucene field type comparison}
\label{tab:fields}
\end{table}

% ─────────────────────────────────────────────────────────────────────────
\section{Retrieval with PyLucene}
% ─────────────────────────────────────────────────────────────────────────

\subsection{IndexSearcher and Similarity}

\begin{lstlisting}[caption={Setting up IndexSearcher with BM25 (from \texttt{src/pylucene/retrieval.py})}]
from org.apache.lucene.index  import DirectoryReader
from org.apache.lucene.search import IndexSearcher
from org.apache.lucene.search.similarities import BM25Similarity

reader   = DirectoryReader.open(directory)
searcher = IndexSearcher(reader)
searcher.setSimilarity(BM25Similarity(1.2, 0.75))  # k1=1.2, b=0.75

from org.apache.lucene.queryparser.classic import QueryParser
parser   = QueryParser("plot", analyzer)
query    = parser.parse("space travel wormhole")
top_docs = searcher.search(query, 5)           # top-5 results

for hit in top_docs.scoreDocs:
    doc = searcher.doc(hit.doc)
    print(f"Score={hit.score:.4f}  |  {doc.get('title')}")
reader.close()
\end{lstlisting}

\subsection{Explaining Scores}

\texttt{IndexSearcher.explain()} returns a human-readable breakdown of how each
relevance score was computed—invaluable for debugging:

\begin{lstlisting}[caption={Score explanation}]
hit         = top_docs.scoreDocs[0]
explanation = searcher.explain(query, hit.doc)
print(explanation.toString())
# Output includes: IDF, TF, field norm, BM25 saturation factor...
\end{lstlisting}

% ─────────────────────────────────────────────────────────────────────────
\section{PyLucene Query Classes}
% ─────────────────────────────────────────────────────────────────────────

PyLucene provides a rich query API. All examples below are drawn from
\texttt{src/pylucene/query\_classes.py} in the repository.

\subsection{TermQuery}

Matches documents containing an \emph{exact} post-analysis term:

\begin{lstlisting}[caption={TermQuery -- exact single-term match}]
from org.apache.lucene.search import TermQuery
from org.apache.lucene.index  import Term

query = TermQuery(Term("plot", "crime"))
top_docs = searcher.search(query, 5)
\end{lstlisting}

\subsection{PhraseQuery}

Ordered phrase match with optional \texttt{slop} for proximity search:

\begin{lstlisting}[caption={PhraseQuery -- exact phrase and proximity}]
from org.apache.lucene.search import PhraseQuery

# Exact phrase
exact = (PhraseQuery.Builder()
         .add(Term("plot", "corporate"))
         .add(Term("plot", "secrets"))
         .build())

# Proximity: terms within 3 positions of each other
prox = (PhraseQuery.Builder()
        .add(Term("plot", "crime"))
        .add(Term("plot", "redemption"))
        .setSlop(3).build())
\end{lstlisting}

\subsection{BooleanQuery}

Combine sub-queries with \texttt{MUST} (AND), \texttt{SHOULD} (OR),
\texttt{MUST\_NOT} (NOT):

\begin{lstlisting}[caption={BooleanQuery -- combining clauses (MUST/SHOULD/MUST NOT)}]
from org.apache.lucene.search import BooleanQuery, BooleanClause

builder = BooleanQuery.Builder()
builder.add(TermQuery(Term("plot",  "crime")),
            BooleanClause.Occur.SHOULD)
builder.add(TermQuery(Term("plot",  "redemption")),
            BooleanClause.Occur.SHOULD)
builder.add(TermQuery(Term("genre", "action")),
            BooleanClause.Occur.MUST_NOT)
query = builder.build()
\end{lstlisting}

\subsection{Numeric Range Queries}

\begin{lstlisting}[caption={IntPoint and FloatPoint -- numeric range queries}]
from org.apache.lucene.document import IntPoint, FloatPoint
from java.lang import Float

# Movies from 1990 to 2000
year_q   = IntPoint.newRangeQuery("year", 1990, 2000)

# Movies with rating >= 9.0
rating_q = FloatPoint.newRangeQuery("rating", 9.0, Float.MAX_VALUE)
\end{lstlisting}

\subsection{WildcardQuery and FuzzyQuery}

\begin{lstlisting}[caption={WildcardQuery and FuzzyQuery -- approximate matching}]
from org.apache.lucene.search import WildcardQuery, FuzzyQuery

# ? = any single char, * = zero or more chars
wild  = WildcardQuery(Term("plot", "r?dempt*"))

# Levenshtein distance <= 2 (handles typos)
fuzzy = FuzzyQuery(Term("plot", "astronawts"), 2)
\end{lstlisting}

\subsection{PrefixQuery and TermRangeQuery}

\begin{lstlisting}[caption={PrefixQuery and TermRangeQuery}]
from org.apache.lucene.search import PrefixQuery, TermRangeQuery

prefix = PrefixQuery(Term("plot", "inter"))           # plot:inter*
trange = TermRangeQuery.newStringRange(
             "genre", "crime", "drama", True, True)   # inclusive
\end{lstlisting}

\subsection{MatchAllDocsQuery and BoostQuery}

\begin{lstlisting}[caption={MatchAllDocsQuery and BoostQuery}]
from org.apache.lucene.search import MatchAllDocsQuery, BoostQuery

all_q   = MatchAllDocsQuery()                        # returns every doc

title_q = TermQuery(Term("title", "redemption"))
boosted = BoostQuery(title_q, 3.0)                   # title match = 3x score
\end{lstlisting}

% ─────────────────────────────────────────────────────────────────────────
\section{Repository Overview}
% ─────────────────────────────────────────────────────────────────────────

The full source code is available at:
\begin{center}
  \href{https://github.com/aayushvij04/IR}{\texttt{https://github.com/aayushvij04/IR}}
\end{center}

\begin{table}[h]
\centering
\small
\begin{tabular}{@{}ll@{}}
\toprule
\textbf{File} & \textbf{Contents} \\
\midrule
\texttt{src/pylucene/indexing.py}     & IndexWriter, Document, Field types, Analyzer choices \\
\texttt{src/pylucene/retrieval.py}    & IndexSearcher, BM25/TF-IDF, multi-field, score explain \\
\texttt{src/pylucene/query\_classes.py} & All query types with runnable demos \\
\texttt{src/boolean\_retrieval.py}    & Inverted index, Boolean ops, skip pointers (pure Python) \\
\texttt{src/ranking\_models.py}       & TF-IDF, VSM, BM25, BM25+ (pure Python) \\
\texttt{src/language\_models.py}      & Unigram LM, JM \& Dirichlet smoothing, KLD \\
\texttt{src/evaluation.py}            & P, R, AP, MAP, nDCG, MRR, Kappa \\
\texttt{src/web\_search.py}           & PageRank, HITS, Shingling \\
\texttt{src/advanced\_ir.py}          & LSI, co-occurrence embeddings, Rocchio \\
\texttt{report/IR\_Report.tex}        & This report (LaTeX source) \\
\bottomrule
\end{tabular}
\caption{Repository file index}
\end{table}

% ─────────────────────────────────────────────────────────────────────────
\section{Conclusion}
% ─────────────────────────────────────────────────────────────────────────

PyLucene brings the full power of Apache Lucene to the Python ecosystem without
sacrificing performance. Its rich field model, analyzer pipeline, and query API
make it suitable for both research prototyping and production search systems.
The integration of BM25 as the default similarity function—combined with
flexible numeric indexing, proximity queries, and fuzzy matching—makes PyLucene
one of the most complete IR toolkits available.

The accompanying code in this repository demonstrates each concept with
runnable, well-commented examples, bridging the gap between IR theory and
practical implementation.

\vspace{1em}
\begin{center}
  \small\textit{Source: \href{https://github.com/aayushvij04/IR}{github.com/aayushvij04/IR}}
\end{center}

\end{document}
